\scanpage{25}

\uuline{$\underline{\overline{\mathrm{XVIII}}}$. $1.2$}

\[
\mathrm I\left\{\begin{aligned}
\mathrm I.a\quad & G\longrightarrow p_*p^*G\\
\mathrm I.b\quad & \underline{\mathrm{Hom}}(G_m,G)\xrightarrow{\sim}{\unsure{R|\mathbb R}}^1p_*G
\end{aligned}\right.
\]
\[
\mathrm{II}\ \left\{\ \text{le foncteur $(1.2.1.1)$ est une équivalence.}\right.
\]

\uline{Remarques 1} : \deleted{\unsure{disparaît}} moyennant $(\mathrm I.a)$ et l'hypothèse d'une section
on a une suite exacte
\[
0\longrightarrow H^1(S,G)\longrightarrow {\unsure{\mathbb H|H}}^{1}(\unsure{\mathbb P|P}(\unsure{\mathscr E|\mathcal E|E}),G)
\xrightarrow{\theta}\Gamma(S,R^1p_*G)\longrightarrow0
\]

\uline{Remarque 2} : La pleine fidélité de $(1.2.1.1)$ est locale sur la base et
moyennant la pleine fidélité, la surjectivité essentielle est locale sur la base

\uline{Remarque 3} : $\mathrm I.b$ (et $\mathrm I.a$) entraîne que
$\unsure{\mathrm{Hom}|\underline{\mathrm{Hom}}}(G_m,G)\longrightarrow H^1(\unsure{\underline{P}|\mathbb P|P}(\mathscr E),G)$ est injectif.

\deleted{$\mathrm I)\Rightarrow\mathrm{II})$ pleine fidélité} :

\inserted{\unsure{on se ramène au cas} où il y a une section.}

$\mathrm I\Rightarrow\mathrm{II}$. \uline{pleine fidélité} :

si $K(k_0,\unsure{p|\varphi})\simeq K(k'_0,\varphi')$ alors $\varphi=\varphi'$
\unsure{Alors} $k_0=k'_0$ et $\mathrm I.a$ entraîne la pleine fidélité.

\uline{essentielle surjectivité} : Soit $k\in H^1(\mathbb P(\unsure{E|\mathcal E}),G)$
$\theta(k)=\varphi(\unsure{O|\mathcal O}(1))$ pour un certain $\varphi$ et
$\theta(k)-\varphi(O(1))\in H^1(S,G)$.

$\mathrm{II})\Rightarrow\mathrm I)$ \quad $\mathrm I\text{-}a)$ résulte de la pleine fidélité.

$\mathrm I.b)$ \quad $\deleted{\mathrm{Hom}(G_m,G)\longrightarrow}$ trivial car a \unsure{quasi-inv$^{\mathrm{se}}$|qui précède}.
% archive sheet number: 24

